\section{Hardware--simulator comparison}
\label{sec:comparison}

Figure~\ref{fig:D08D09}(c,d) compares the combined quantity $Q$ from simulations at $p=1,2$ and hardware at $p=1$. Measurements on the Baihua processor are organized into repeated acquisition units. Each unit contains $15$ values of $s$ along two scans, $\delta\in\{0.85,-2\}$, giving $30$ parameter points with $1024$ shots per point and one additional pair of readout-calibration circuits. A single round of $500$ units is analyzed. Error bars on hardware points indicate standard errors across acquisition units.

The observables $\hat S(\pi)$ and $\hat O_{\rm str}$ are evaluated from the same readout-mitigated distributions. Hardware hats keep the same sign convention ($\hat O=-O_{\rm str}$) and are divided by the selected window's shared maxima: the top-1 maxima of the window-pooled $30$-point means, refitted for each candidate window. For the selected window these are $S_{\max}=4.51$ and $O_{\max}=0.090$. Simulations use fixed scales, $\hat S=S/8$ and $\hat O=-O_{\rm str}$, and the $p=1$ simulation is the reference curve for window selection ($p=2$ is shown for comparison only). Both use
\begin{equation}
Q=\frac{4}{3}-2\hat O-\frac{4}{3}\hat S.
\end{equation}
With the fixed simulation scales, this expression agrees with the definition in Sec.~\ref{sec:heatmaps}. The combined $Q$ values are used directly, with no post-combination renormalization, to evaluate the distance metric $d_{\rm raw}$, the mean absolute $Q$ distance over the $30$ points, and the Pearson correlation coefficient $r$.

Among the compared even-length windows, the best window was $[226,242)$, comprising $16$ acquisition units, with $d_{\rm raw}=0.30$ and $r=0.85$. Candidate windows are all contiguous even-length subsegments (length at least $10$) of the $500$-unit sequence, each with its own refitted endpoints; the best window maximizes $r$ against the $p=1$ simulation, with ties broken by smaller $d_{\rm raw}$, then longer length, then earlier start. These metrics are report-only descriptions of curve agreement under the selected window and its normalization.

Along $\delta=-2$, hardware results exhibit a parameter dependence similar to the simulated $p=2$ curve throughout the scan. At $\delta=0.85$, the first half follows a similar trend, while in the topological region hardware $Q$ turns negative and reaches lower values than either simulation as $s\to1$. A resolved string signal occurs at $s=0.99$: pooling all $N=500$ units, the hardware $\hat O_{\rm str}$ mean is approximately $5\%$ of the theoretical value, $9\sigma$ away from zero where $\sigma$ is the inter-unit standard error; the remaining points are consistent with zero within measurement uncertainty.
